\documentclass[11pt,a4paper]{article}
\usepackage[margin=2.5cm,headheight=15pt]{geometry}
\usepackage{amsmath,amssymb,amsthm}
\usepackage{mathtools}
\usepackage{microtype}
\usepackage{hyperref}
\usepackage{cleveref}
\usepackage{enumitem}
\usepackage{mdframed}
\usepackage{xcolor}
\usepackage{titlesec}
\usepackage{booktabs}
\usepackage{fancyhdr}

\definecolor{navyblue}{RGB}{26,58,92}
\definecolor{midblue}{RGB}{44,95,138}
\definecolor{lightblue}{RGB}{232,240,248}
\definecolor{verylightblue}{RGB}{248,251,254}
\definecolor{proofgray}{RGB}{242,242,242}
\definecolor{warncolor}{RGB}{255,245,230}
\definecolor{warnborder}{RGB}{200,100,0}
\definecolor{opencolor}{RGB}{235,245,235}
\definecolor{openborder}{RGB}{0,120,0}

\hypersetup{colorlinks=true,linkcolor=midblue,citecolor=midblue,urlcolor=midblue,
  pdftitle={SSP Part VIII: Branch-and-Benders}}

\pagestyle{fancy}\fancyhf{}
\fancyhead[L]{\small\color{navyblue}\textit{Part VIII: Branch-and-Benders for the Uniform SSP}}
\fancyhead[R]{\small\color{navyblue}\thepage}
\renewcommand{\headrulewidth}{0.4pt}

\titleformat{\section}{\large\bfseries\color{navyblue}}{}{0em}{}[\color{navyblue}\titlerule]
\titleformat{\subsection}{\normalsize\bfseries\color{midblue}}{}{0em}{}

\newmdenv[backgroundcolor=verylightblue,linecolor=navyblue,linewidth=1.2pt,
  leftline=true,rightline=false,topline=false,bottomline=false,
  innerleftmargin=12pt,innerrightmargin=8pt,innertopmargin=8pt,innerbottommargin=8pt,
  skipabove=10pt,skipbelow=6pt]{thmbox}
\newmdenv[backgroundcolor=proofgray,linecolor=midblue,linewidth=0.8pt,
  leftline=true,rightline=false,topline=false,bottomline=false,
  innerleftmargin=10pt,innerrightmargin=8pt,innertopmargin=6pt,innerbottommargin=6pt,
  skipabove=6pt,skipbelow=6pt]{proofbox}
\newmdenv[backgroundcolor=lightblue,linecolor=navyblue,linewidth=1.5pt,roundcorner=3pt,
  innerleftmargin=14pt,innerrightmargin=14pt,innertopmargin=10pt,innerbottommargin=10pt,
  skipabove=12pt,skipbelow=8pt]{formbox}
\newmdenv[backgroundcolor=warncolor,linecolor=warnborder,linewidth=1.2pt,
  leftline=true,rightline=false,topline=false,bottomline=false,
  innerleftmargin=12pt,innerrightmargin=8pt,innertopmargin=8pt,innerbottommargin=8pt,
  skipabove=10pt,skipbelow=6pt]{warnbox}
\newmdenv[backgroundcolor=opencolor,linecolor=openborder,linewidth=1.2pt,
  leftline=true,rightline=false,topline=false,bottomline=false,
  innerleftmargin=12pt,innerrightmargin=8pt,innertopmargin=8pt,innerbottommargin=8pt,
  skipabove=10pt,skipbelow=6pt]{openproblem}

\theoremstyle{plain}
\newtheorem{theorem}{Theorem}[section]
\newtheorem{lemma}[theorem]{Lemma}
\newtheorem{proposition}[theorem]{Proposition}
\theoremstyle{definition}
\newtheorem{definition}[theorem]{Definition}
\newtheorem{remark}[theorem]{Remark}

\DeclareMathOperator{\KTNS}{KTNS}
\DeclareMathOperator{\OPT}{OPT}
\newcommand{\ZSSP}{Z_{\mathrm{SSP}}^{*}}
\newcommand{\Kstar}{K^{*}}
\setlist[itemize]{noitemsep,topsep=4pt}
\setlist[enumerate]{noitemsep,topsep=4pt}

\begin{document}

\begin{center}
  {\color{navyblue}\rule{\textwidth}{2pt}}\\[0.4cm]
  {\LARGE\bfseries\color{navyblue}The Job Sequencing and Tool Switching Problem}\\[0.2cm]
  {\large\bfseries\color{navyblue}Part VIII: Branch-and-Benders Decomposition (BBC)}\\[0.1cm]
  {\large\itshape\color{midblue}A master sequence problem with a polynomial KTNS subproblem}\\[0.3cm]
  {\color{navyblue}\rule{\textwidth}{2pt}}
\end{center}

\vspace{0.3cm}
\begin{mdframed}[backgroundcolor=lightblue,linecolor=navyblue,linewidth=1pt,
  innerleftmargin=14pt,innerrightmargin=14pt,innertopmargin=8pt,innerbottommargin=8pt]
\small\textbf{Scope and status.}
This note develops the Branch-and-Benders-cut (BBC) exact method for the uniform SSP. The
SSP splits cleanly into a \emph{sequencing} decision and a \emph{tooling} decision: fix the
job sequence and the optimal tooling cost is computed in polynomial time by KTNS
\cite{Tang1988}. BBC puts the sequence in a \emph{master} ATSP (with a surrogate
$\theta$ for the tool cost) and the tooling in a \emph{subproblem} solved by KTNS, linked by
Benders optimality cuts. The decisive structural feature --- and the reason BBC is
attractive here --- is that \textbf{the subproblem is polynomial}, in contrast to the
Dantzig--Wolfe/column-generation route of Part~II whose pricing subproblem is itself an
NP-hard SSP. The integer (lazy) version is implemented and exact
(\texttt{src/BBC/branch\_and\_benders\_cut\_cplex.py}); separating Benders cuts at
\emph{fractional} master solutions is an open direction (\cref{sec:open}).
\end{mdframed}

\vspace{0.3cm}
\tableofcontents
\newpage

% -----------------------------------------------------------------------------
\section{The Decomposition}
\label{sec:decomp}
% -----------------------------------------------------------------------------

\paragraph{Notation.} Jobs $J=\{1,\dots,n\}$, tools $T$, capacity $b$, required sets
$T_j\subseteq T$. A dummy job $0$ with $T_0=\emptyset$ acts as a free start/end depot, so the
schedule is an open path over the real jobs (no tool cost is charged for loading the first
configuration). For a job sequence $\sigma$, $\KTNS(\sigma)$ denotes the minimum number of
tool switches over $\sigma$, computed by the Keep-Tool-Needed-Soonest policy.

\paragraph{The split.} A feasible SSP solution is a pair (sequence $\sigma$, tooling). For a
\emph{fixed} $\sigma$ the optimal tooling is decided independently and exactly:

\begin{thmbox}
\begin{theorem}[KTNS solves the subproblem exactly and in polynomial time \cite{Tang1988}]
\label{thm:ktns}
For any fixed job sequence $\sigma$, the minimum number of tool switches equals
$\KTNS(\sigma)$, computable in $O(n\,|T|)$ time.
\end{theorem}
\end{thmbox}

Hence
\[
  \ZSSP \;=\; \min_{\sigma}\ \KTNS(\sigma),
\]
the minimisation being over the $n!$ job sequences (a permutation problem). BBC optimises
this by branch-and-cut on a master that chooses $\sigma$, deferring the evaluation
$\KTNS(\sigma)$ to a callback.

\begin{remark}[Why this split is the right one for the SSP]
Two decompositions compete (cf.\ Parts~II and~IV). \emph{Dantzig--Wolfe / column generation}
(Part~II) generates configurations as columns; its pricing subproblem is a set-union knapsack
and, for whole sequences, an NP-hard SSP. \emph{Benders} fixes the complicating variables
(the sequence) and leaves a subproblem that here is \emph{polynomial} (KTNS). Benders
therefore guarantees that exact cuts can always be produced cheaply --- the property the CG
route lacks. The price is that the master's lower bound is only as strong as the cuts
accumulated so far.
\end{remark}

% -----------------------------------------------------------------------------
\section{The Master Problem}
\label{sec:master}
% -----------------------------------------------------------------------------

The master is an asymmetric TSP over $J\cup\{0\}$ with arc variables $x_{ij}\in\{0,1\}$
($x_{ij}=1$ iff job $j$ immediately follows job $i$) and a single continuous surrogate
$\theta\ge 0$ standing for the tool-switching cost of the chosen sequence.

\begin{formbox}
\textbf{BBC master (relaxed; cuts added dynamically)}
\begin{align}
  \min\ \ & \theta \label{eq:m_obj}\\
  \text{s.t.}\ \ & \sum_{j\neq i} x_{ij}=1,\quad \sum_{i\neq j} x_{ij}=1
     && \forall\, i,j\in J\cup\{0\} \label{eq:m_deg}\\
  & \text{(subtour elimination on } x) && \text{(lazy SEC or MTZ, \cref{rem:subtour})}\label{eq:m_sec}\\
  & \theta \ \ge\ \sum_{i\neq j} w_{ij}\, x_{ij},\qquad w_{ij}=\max\!\bigl(0,\ |T_i\cup T_j|-b\bigr)
     \label{eq:m_pair}\\
  & \theta \ \ge\ \text{(Benders optimality cuts, \cref{sec:cuts})} \label{eq:m_benders}\\
  & x_{ij}\in\{0,1\},\quad \theta\ge 0. \notag
\end{align}
\end{formbox}

Constraints \eqref{eq:m_deg}--\eqref{eq:m_sec} make $x$ a Hamiltonian cycle through the
depot, i.e.\ an open job sequence $\sigma(x)$. Inequality \eqref{eq:m_pair} is the
\emph{pairwise lower bound}: each executed transition $i\to j$ forces at least
$w_{ij}=\max(0,|T_i\cup T_j|-b)$ switches, because the two jobs' tools cannot both fit unless
$|T_i\cup T_j|\le b$ \cite{Tang1988,Laporte2004}. It is valid \emph{a priori} and seeds the
master with a non-trivial bound before any Benders cut is generated.

\begin{remark}[Subtour elimination in the master]
\label{rem:subtour}
Any TSP subtour scheme works for \eqref{eq:m_sec}: lazy DFJ cuts (separated on integer
incumbents) or the polynomial MTZ constraints of Part~IV
(\texttt{04\_mtz\_formulation.tex}). Note the roles differ from Part~IV: there MTZ replaced
\emph{configuration}-graph subtour constraints in a non-compact model; here it would only
order the $n$ \emph{jobs}, where it is unproblematic and exact.
\end{remark}

% -----------------------------------------------------------------------------
\section{Benders Optimality Cuts}
\label{sec:cuts}
% -----------------------------------------------------------------------------

The subproblem has no linear-programming dual (KTNS is a combinatorial algorithm), so the
cuts are \emph{logic-based} Benders cuts \cite{Vanderbeck2010}. Given an integer master
incumbent with sequence $\sigma$, arc set $A(\sigma)=\{(i,j):x_{ij}=1\}$, evaluate
$q=\KTNS(\sigma)$. If $\theta<q$ the incumbent is infeasible for the true cost and we add:

\begin{thmbox}
\begin{proposition}[Validity of the combinatorial optimality cut]
\label{prop:cut}
For the sequence $\sigma$ with $q=\KTNS(\sigma)$ and $|A(\sigma)|=n$ (a depot cycle), the
inequality
\[
  \theta \ \ge\ q\Bigl(1-\!\!\sum_{(i,j)\in A(\sigma)}\!\!\bigl(1-x_{ij}\bigr)\Bigr)
  \;=\; q\Bigl(\textstyle\sum_{(i,j)\in A(\sigma)} x_{ij}-(n-1)\Bigr)
\]
is valid: it forces $\theta\ge q$ whenever the master reproduces $\sigma$ exactly, and is
slack ($\le 0$ on the right) for every other sequence.
\end{proposition}
\end{thmbox}
\begin{proofbox}
\begin{proof}
If $x$ reproduces $\sigma$ then $\sum_{(i,j)\in A(\sigma)}x_{ij}=n$, the right-hand side is
$q$, and $\theta\ge q=\KTNS(\sigma)$ is exactly the subproblem optimum (\cref{thm:ktns}),
hence valid. If $x$ omits at least one arc of $\sigma$ then
$\sum_{(i,j)\in A(\sigma)}x_{ij}\le n-1$, the right-hand side is $\le 0\le\theta$, so the cut
imposes nothing. No feasible sequence is wrongly excluded.
\end{proof}
\end{proofbox}

\begin{warnbox}
\textbf{These ``no-good'' cuts are valid but weak.} Each cut binds only on the single
sequence that generated it, so naively BBC could enumerate exponentially many sequences. Two
strengthenings are used/possible: (i)~the always-on pairwise bound \eqref{eq:m_pair} keeps
the root bound non-trivial; (ii)~\emph{cut lifting} --- because $\KTNS(\sigma)$ depends only
on the relative order in which tools are first/last needed (its $0$-block structure, cf.\
Part~I), a single evaluation certifies the same bound for the whole family of sequences
sharing that structure, yielding cuts of the form $\theta\ge q\,(\sum_{(i,j)\in S}x_{ij}-|S|+1)$
for a \emph{subset} $S\subseteq A(\sigma)$ of order-critical arcs. Identifying minimal such
$S$ is the practical lever for BBC performance.
\end{warnbox}

% -----------------------------------------------------------------------------
\section{The Algorithm}
\label{sec:algo}
% -----------------------------------------------------------------------------

BBC is a single branch-and-bound tree on the master with a lazy-constraint callback:

\begin{thmbox}
\textbf{Branch-and-Benders-cut for the uniform SSP (integer cuts).}
\begin{enumerate}[label=\textbf{\arabic*.},leftmargin=2.2em]
\item Build the master \eqref{eq:m_obj}--\eqref{eq:m_pair}; install a lazy-constraint callback.
\item At each integer incumbent $(x,\theta)$:
\begin{enumerate}[label=(\alph*),leftmargin=2em]
\item extract the sequence $\sigma\gets\sigma(x)$; if $x$ contains a subtour, add a subtour
  cut and reject;
\item compute $q\gets \KTNS(\sigma)$ \ ($O(n|T|)$, exact by \cref{thm:ktns});
\item if $\theta<q$, add the optimality cut of \cref{prop:cut} (lifted if possible) and
  reject the incumbent;
\item otherwise accept it --- a genuine optimal-tooling schedule of cost $q=\theta$.
\end{enumerate}
\item Continue until the tree is exhausted; return the incumbent, of cost $\ZSSP$.
\end{enumerate}
\end{thmbox}

\begin{thmbox}
\begin{theorem}[Exactness and finiteness]
\label{thm:exact}
The integer-cut BBC of \cref{sec:algo} terminates and returns $\ZSSP$.
\end{theorem}
\end{thmbox}
\begin{proofbox}
\begin{proof}
Every accepted incumbent satisfies $\theta=q=\KTNS(\sigma)$, hence corresponds to a real
schedule of cost $\theta$; thus the returned value is $\ge\ZSSP$. Conversely the cuts of
\cref{prop:cut} never remove a sequence's true cost (they are valid, \cref{prop:cut}), so the
master remains a relaxation and its optimum is $\le\ZSSP$. Finiteness: there are finitely
many sequences; each rejected incumbent adds a cut that forbids re-accepting that sequence
with too small a $\theta$, so no sequence is rejected twice for the same reason.
\end{proof}
\end{proofbox}

% -----------------------------------------------------------------------------
\section{Position Among the Exact Methods}
\label{sec:position}
% -----------------------------------------------------------------------------

\begin{center}
\renewcommand{\arraystretch}{1.3}
\begin{tabular}{lccc}
\toprule
\textbf{Method} & \textbf{Subproblem} & \textbf{Subproblem cost} & \textbf{Exact?} \\
\midrule
Compact ILP (Part I) & --- & --- & yes (weak LP) \\
DW / column generation (Part II) & pricing (set-union knapsack / SSP) & \textbf{NP-hard} & yes \\
Cluster-MTZ (Part IV) & --- & --- & \emph{conditional} (exchange arg.) \\
\textbf{Branch-and-Benders (this Part)} & tooling via KTNS & \textbf{poly $O(n|T|)$} & yes \\
\bottomrule
\end{tabular}
\end{center}

BBC is the only route here whose subproblem is provably polynomial and whose exactness is
unconditional. Its weakness is dual: the logic-based cuts can be weak, so the master bound
may improve slowly without lifting. This is the mirror image of Part~IV (cheap polynomial
constraints, but a conditional/heuristic model) and of Part~II (a strong LP bound, but an
NP-hard subproblem).

% -----------------------------------------------------------------------------
\section{Open Direction: Fractional Benders Cuts}
\label{sec:open}
% -----------------------------------------------------------------------------

\begin{openproblem}
\textbf{OP (fractional Benders cuts; cf.\ \texttt{09\_open\_problems.tex}, OP\,8).} The implemented BBC separates cuts
only at \emph{integer} incumbents, so the LP relaxation at a node is bounded only by the
pairwise inequality \eqref{eq:m_pair} plus previously added integer cuts. Develop a
separation that produces a \emph{valid} cut at a \emph{fractional} master point
$\bar x$ --- e.g.\ from a relaxation of the KTNS recursion (a min-cost-flow / interval-matrix
lower bound on the tooling cost as a function of fractional arc usage), or from the totally
unimodular tooling LP of Catanzaro et al.\ for fixed-support sequences. Such user cuts would
tighten the bound \emph{before} branching and are the most promising avenue for scaling BBC
beyond the sizes the integer-cut version reaches. Whether a polynomially separable family of
fractional Benders cuts exists for the KTNS subproblem is open.
\end{openproblem}

\begin{remark}[Relation to the gap analysis]
By the grouping-exactness theorem of Part~V, $\ZSSP=\min_{\mathcal{G}}\mathrm{GTSP}(\mathcal{G})$;
BBC reaches the same optimum from the sequence side. A primal heuristic for the BBC incumbent
is therefore the JGP$+$GSP schedule of Part~V, and the bound
$\theta\ge\max(\Kstar-1,\,|U|-b)$ (Part~V) may be added to \eqref{eq:m_pair} as an extra valid
inequality to strengthen the root.
\end{remark}

% -----------------------------------------------------------------------------
\section{Implementation Notes}
\label{sec:impl}
% -----------------------------------------------------------------------------

The CPLEX implementation is \texttt{src/BBC/branch\_and\_benders\_cut\_cplex.py} (public API
in \texttt{branch\_and\_benders\_cut.py}; shared subtour/sequence/KTNS utilities in
\texttt{bbc\_common.py}). The KTNS evaluation and subtour detection are in the lazy callback;
the dummy depot gives the free initial load, so the reported objective is the number of
inter-job switches (consistent with the convention of Parts~I and~V). Comparison baselines
(LSS, SSPMF) are in \texttt{lss\_formulation.py} and \texttt{sspmf\_formulation.py}.

\bibliographystyle{abbrvnat}
\bibliography{references}
\end{document}
